% ============================================================
% supplementary.tex — 补充材料（Stage 1B 重写版）
% 对应正文：main_cn.tex（Stage 1B 表型原型）
%
% 编译方法：
%   xelatex supplementary.tex
%
% 注意：此文件独立编译，不依赖 bibtex
% ============================================================

\documentclass[12pt, a4paper]{ctexart}

% ------ 基础包 ------
\usepackage{geometry}
\geometry{left=2.5cm, right=2.5cm, top=3cm, bottom=3cm}
\usepackage{setspace}
\onehalfspacing

% ------ 参考文献 ------
\usepackage[super,square,sort&compress]{natbib}
\usepackage{gbt7714}
\bibliographystyle{gbt7714-numerical}

% ------ 图表 ------
\usepackage{graphicx}
\usepackage{booktabs}
\usepackage{array}
\usepackage{multirow}
\usepackage{caption}
\captionsetup{font=small, labelfont=bf}

% ------ 数学 ------
\usepackage{amsmath}
\usepackage{amssymb}

% ------ 颜色与链接 ------
\usepackage[dvipsnames]{xcolor}
\usepackage[hidelinks]{hyperref}

% ------ 其他 ------
\usepackage{enumitem}
\usepackage{float}

% 补充材料编号重置：表格前缀 S，图片前缀 S
\renewcommand{\thetable}{S\arabic{table}}
\renewcommand{\thefigure}{S\arabic{figure}}

% ------ 元数据 ------
\title{补充材料\\
\large 基于汇总级心肺运动试验的运动安全实验室表型原型构建}

\author{\textcolor{red}{TODO: 作者姓名}}
\date{2026年4月}

% ============================================================
\begin{document}
\maketitle

\tableofcontents
\clearpage

% ============================================================
\section*{补充方法}
\addcontentsline{toc}{section}{补充方法}
% ============================================================

\subsection*{S1. 变量角色分类详细定义}
\addcontentsline{toc}{subsection}{S1. 变量角色分类详细定义}

Stage 1B 框架中所有 CPET 变量严格按以下四类分配，防止构念污染：

\begin{table}[htbp]
\centering
\small
\caption{Stage 1B 变量角色分类}
\label{tab:s1_variable_roles}
\begin{tabular}{llll}
\toprule
角色类型 & 变量名 & 域/层 & 用途 \\
\midrule
表型变量（Reserve 域）& vo2\_peak & Reserve & 负担评分 \\
 & vt1\_vo2 & Reserve & 负担评分 \\
 & o2\_pulse\_peak & Reserve & 负担评分 \\
 & mets\_peak & Reserve & 负担评分 \\
\midrule
表型变量（Ventilatory 域）& ve\_vco2\_slope & Ventilatory & 负担评分 \\
 & oues & Ventilatory & 负担评分 \\
\midrule
不稳定变量 & eih\_status & Instability & 覆盖规则（Severe） \\
 & bp\_peak\_sys & Instability & 覆盖规则（Severe/Mild） \\
 & bp\_peak\_dia & Instability & 覆盖规则（Severe） \\
 & o2\_pulse\_trajectory & Instability & 覆盖规则（Severe） \\
\midrule
置信度辅助变量 & hr\_peak\_pct\_pred & Effort & 努力程度代理 \\
 & vo2\_peak\_pct\_pred & Anchor & 锚点一致性 \\
\midrule
验证变量 & test\_result & Validation & 外部锚定验证（仅） \\
 & group\_code & Validation & 队列编码 \\
\midrule
排除变量 & zone\_v2, p1\_zone & Legacy & 仅 Legacy 对照 \\
 & S\_lab & Legacy & 仅 Legacy 对照 \\
\bottomrule
\end{tabular}
\\[1ex]
{\footnotesize test\_result 在 final\_zone 确定后方可使用，严禁用于定义区间边界。}
\end{table}

\subsection*{S2. 条件分位参考模型技术细节}
\addcontentsline{toc}{subsection}{S2. 条件分位参考模型技术细节}

\textbf{模型架构}：sklearn QuantileRegressor（线性分位数回归）。
特征工程：SplineTransformer（knots=5，degree=3）处理年龄连续效应；
OneHotEncoder（handle\_unknown='ignore'）处理性别（男/女）和方案（踏车/平板）。
正则化：alpha=0.001（默认），防止小样本 strict reference 过拟合。

\textbf{单调性约束}：对每个样本，5个分位数预测值经 np.sort 升序排列，
保证 q10 $\leq$ q25 $\leq$ q50 $\leq$ q75 $\leq$ q90 严格成立。

\textbf{分位模型评估}：在 wide reference subset 上验证，
check: q10 实际覆盖率 $\approx$10\%，q90 实际覆盖率 $\approx$90\%
（分位数校准检验）。

\textbf{目标变量}：6个核心表型变量分别独立拟合分位模型，
各变量的分位束（QuantileModelBundle）保存为 joblib 格式（
\texttt{outputs/reference\_models/quantile\_bundle\_stage1b.joblib}）。

\subsection*{S3. 表型负担评分技术细节}
\addcontentsline{toc}{subsection}{S3. 表型负担评分技术细节}

\textbf{变量级负担转换规则}（表~\ref{tab:s2_burden_rules}）：

\begin{table}[htbp]
\centering
\small
\caption{表型负担转换规则}
\label{tab:s2_burden_rules}
\begin{tabular}{lllll}
\toprule
变量 & 方向 & 负担=0（正常）& 负担=0.5（临界）& 负担=1.0（异常）\\
\midrule
vo2\_peak & higher-better & $\geq$q25 & [q10, q25) & $<$q10 \\
vt1\_vo2 & higher-better & $\geq$q25 & [q10, q25) & $<$q10 \\
o2\_pulse\_peak & higher-better & $\geq$q25 & [q10, q25) & $<$q10 \\
mets\_peak & higher-better & $\geq$q25 & [q10, q25) & $<$q10 \\
ve\_vco2\_slope & higher-worse & $\leq$q75 & (q75, q90] & $>$q90 \\
oues & higher-better & $\geq$q25 & [q10, q25) & $<$q10 \\
\bottomrule
\end{tabular}
\end{table}

\textbf{缺失值处理}：字段缺失时该变量负担记 NaN，
域内均值以可用字段平均（忽略 NaN）；
若 Reserve 域可用字段 $<$2 或 Ventilatory 域可用字段 $<$1，
则对应域评分记 NaN，P\_lab 可能为 NaN。

\textbf{表型区间切点}：
以 strict reference subset 上 P\_lab 分布的 P75 和 P90 为切点，
两点均由参考子集自动估计；若切点退化（P75=P90），
则 P90 替换为 P75+epsilon（0.001）以避免空区间。

\subsection*{S4. 置信度权重来源}
\addcontentsline{toc}{subsection}{S4. 置信度权重来源}

置信度引擎四分量权重（0.40/0.15/0.20/0.25）基于以下方法学先验：

\begin{table}[htbp]
\centering
\small
\caption{置信度分量权重说明}
\label{tab:s3_confidence_weights}
\begin{tabular}{llcl}
\toprule
分量 & 含义 & 权重 & 理由 \\
\midrule
$C_{\text{completeness}}$ & 核心字段完整率 & 0.40 & 数据质量是置信度的最主要来源 \\
$C_{\text{effort}}$ & 努力程度（HR\%pred） & 0.15 & 间接代理，信息量有限 \\
$C_{\text{anchor}}$ & 外部-内部参考一致 & 0.20 & 参考锚点验证 \\
$C_{\text{validation}}$ & 与结局三分位一致 & 0.25 & 结局方向验证，但依赖结局模型质量 \\
\bottomrule
\end{tabular}
\\[1ex]
{\footnotesize 权重来自方法学先验，未经前瞻性数据校准。
未来研究应以标注队列（含真实结局和随访数据）校准各分量权重。}
\end{table}

\subsection*{S5. Stage 1B 验收标准}
\addcontentsline{toc}{subsection}{S5. Stage 1B 验收标准}

管线输出的验收判定基于以下标准（\texttt{stage1b\_report.py} 实现）：

\begin{table}[htbp]
\centering
\small
\caption{Stage 1B 验收标准与判定逻辑}
\label{tab:s4_acceptance}
\begin{tabular}{llll}
\toprule
指标 & 阈值 & 判定级别 & 说明 \\
\midrule
strict reference 中 red\% & $<$ 15\% & Fail if violated & 参考合理性 \\
strict reference 中 green\% & $>$ 50\% & Warn if violated & 参考合理性 \\
final\_zone vs test\_result 梯度 & 方向正确（Green$<$Yellow$<$Red）& Fail if reversed & 构念效度 \\
 & 非单调但非反向 & Warn if non-monotone & 构念效度 \\
high confidence 比例 & $\geq$ 10\% & Warn if below & 置信度分布 \\
\midrule
\multicolumn{4}{l}{\textit{综合判定：存在 Fail 条件 → Fail；仅 Warn → Warn；否则 → Accept}} \\
\bottomrule
\end{tabular}
\end{table}

\clearpage

% ============================================================
\section*{补充结果}
\addcontentsline{toc}{section}{补充结果}
% ============================================================

\subsection*{S6. Stage 1B 管线输出表结构}
\addcontentsline{toc}{subsection}{S6. Stage 1B 管线输出表结构}

\texttt{make stage1b} 生成的 \texttt{data/output/stage1b\_output\_table.parquet}
包含以下 13+ 列（每人一行）：

\begin{table}[htbp]
\centering
\small
\caption{Stage 1B 输出表列定义}
\label{tab:s5_output_columns}
\begin{tabular}{lll}
\toprule
列名 & 类型 & 含义 \\
\midrule
reserve\_burden & float [0,1] & Reserve 域平均负担 \\
vent\_burden & float [0,1] & Ventilatory 域平均负担 \\
p\_lab & float [0,1] & 综合表型负担（0.5×R + 0.5×V） \\
phenotype\_zone & str & 表型区（green/yellow/red） \\
instability\_severe & bool & 严重不稳定标志 \\
instability\_mild & bool & 轻度不稳定标志 \\
final\_zone\_before\_confidence & str & 覆盖后但置信前区间 \\
confidence\_score & float [0,1] & 综合置信度 \\
confidence\_label & str & high/medium/low \\
indeterminate\_flag & bool & 不确定区标志（True=yellow\_gray） \\
final\_zone & str & 最终区间（含 yellow\_gray） \\
completeness\_score & float [0,1] & 字段完整率 \\
effort\_score & float [0,1] & 努力程度分 \\
outcome\_risk\_prob & float [0,1] & 结局锚点模型预测概率（可选） \\
outcome\_risk\_tertile & str & low/mid/high（可选） \\
anomaly\_score & float & Mahalanobis 距离（QC用，可选） \\
anomaly\_flag & bool & 异常表型标志（QC用，可选） \\
\bottomrule
\end{tabular}
\end{table}

\subsection*{S7. Legacy 系统对照（Phase G vs Stage 1B）}
\addcontentsline{toc}{subsection}{S7. Legacy 系统对照（Phase G vs Stage 1B）}

\begin{table}[htbp]
\centering
\small
\caption{Phase G S\_lab 系统 vs Stage 1B 表型原型对照}
\label{tab:s6_legacy_compare}
\begin{tabular}{lll}
\toprule
维度 & Phase G（S\_lab，legacy）& Stage 1B（表型原型） \\
\midrule
核心评分 & 等权 R/T/I 三轴连续评分（0--100）& 表型负担评分（P\_lab，0--1） \\
参考方程 & OLS（年龄+性别+BMI）& 条件分位（年龄样条+性别+方案） \\
不稳定处理 & 混合于 I 轴连续评分 & 独立覆盖规则（Severe/Mild） \\
test\_result 角色 & 外部验证 & 外部验证（相同，但更清晰的边界定义） \\
不确定性输出 & 无（强制三分类）& 显式 yellow\_gray（indeterminate）\\
置信度 & 无 & 四分量加权置信度（0.40/0.15/0.20/0.25）\\
Legacy ML & P0/P1 ML 预测（AUC 0.58，F1 0.47--0.50）& 作为 Legacy 对照保留 \\
异常审计 & Phase G 独立 Mahalanobis 模型（方向不一致）& Robust MinCovDet，定位 QC \\
验收标准 & 无明确标准 & Accept/Warn/Fail（3项量化标准）\\
\bottomrule
\end{tabular}
\\[1ex]
{\footnotesize Stage 1B 框架设计的关键改进：构念分离、参考精细化、不确定性显式化。
Phase G 结果（P0 AUC=0.583，P1 F1=0.499，结局模型 AUC=0.548）
保存于 \texttt{reports/legacy/} 和 \texttt{reports/legacy\_baseline\_manifest.json}，
供 legacy 对照使用，不参与 Stage 1B 主要结论。}
\end{table}

\subsection*{S8. Stage 1B 输出结果（待填充）}
\addcontentsline{toc}{subsection}{S8. Stage 1B 输出结果（待填充）}

以下各项已于 \texttt{make stage1b} 运行于真实数据（N=3232）后填充，
来源：\texttt{reports/stage1b\_summary\_report.md}（2026-04-22）。

\begin{table}[htbp]
\centering
\small
\caption{Stage 1B 全量输出摘要（N=3232，2026-04-22 真实数据运行结果）}
\label{tab:s7_output_summary}
\begin{tabular}{lcc}
\toprule
指标 & 值 & 说明 \\
\midrule
总样本数 & 3,232 & 全量 staging parquet \\
strict reference 子集 N & 855（26.5\%）& CTRL + 严格标准 \\
phenotype Green\% & 30.7\%（992例）& 全样本 \\
phenotype Yellow\% & 27.1\%（877例）& 全样本 \\
phenotype Red\% & 41.0\%（1326例）& 全样本 \\
Severe instability N（\%）& 220（6.8\%）& 强制 Red 样本 \\
Mild instability 升级 N（\%）& 243（7.5\%）& Green→Yellow 覆盖变更 \\
final\_zone Green\% & 25.0\%（807例）& 含覆盖+置信修正 \\
final\_zone Yellow\% & 28.9\%（933例）& \\
final\_zone Red\% & 44.3\%（1432例）& \\
yellow\_gray（indeterminate）\% & 1.9\%（60例）& \\
High confidence\% & 74.2\%（2397例）& $\geq$0.75 \\
test\_result 梯度方向 & non-monotone & Green 14.1\%，Yellow 15.2\%，Red 14.7\% \\
结局模型 CV AUC & 0.564（±0.016）& 预期 0.55--0.65，Test AUC=0.608 \\
验收判定 & \textbf{Warn} & 参考标准通过；梯度不完整 \\
\bottomrule
\end{tabular}
\end{table}

\subsection*{S9. 旧 Phase G ML 辅助验证结果（Legacy 保留）}
\addcontentsline{toc}{subsection}{S9. 旧 Phase G ML 辅助验证结果（Legacy 保留）}

以下结果来自 Phase G（legacy），保留供方法学对照参考，
不参与 Stage 1B 主要结论。

\textbf{运动前风险预测（P0）}：
XGBoost（无BP特征）AUC=0.583，CV-AUC=0.623$\pm$0.028；
2变量基线（HTN+BMI）AUC=0.611 $>$ ML 全特征（AUC=0.583），
反映汇总级数据的预测性能上限。

\textbf{运动后安全区分类（P1）}：
CatBoost F1$_\text{macro}$=0.499，$\kappa$=0.302；
Green 召回67.1\%，Yellow 召回51.8\%，Red 召回30.8\%。
P1 代理泄漏消融分析：Full vs No-VO2 $\Delta$F1=+0.003（可忽略），
证实预测性能上限来自特征集固有局限，而非泄漏驱动。

\textbf{结局锚定模型（Phase G）}：
CV AUC=0.578$\pm$0.031，测试集 AUC=0.548；
Red 区阳性率（25.3\%）$>$ Yellow 区（10.5\%）$>$ Green 区（5.1\%）呈单调梯度。
此结果在 Stage 1B 框架中复现于结局锚点验证器，
验证了 final\_zone 对 test\_result 的构念效度方向。

\textbf{Mahalanobis 异常区（Phase G）}：
经验切点 P75=3.825，P95=8.788；参考子集 N=969；
$D^2$ vs test\_result 相关 $r$=$-$0.011（方向性弱）。
Stage 1B 中重新定位为 QC 审计（anomaly\_flag），
不再参与 zone 定义。

\subsection*{S10. 多定义一致性框架（Phase G，Legacy）}
\addcontentsline{toc}{subsection}{S10. 多定义一致性框架（Phase G，Legacy）}

Phase G 以 K=4 定义投票分析多定义一致性：
高信度841例（26.2\%），不确定区2,358例（73.5\%），Green/Red 冲突1,974例（61.6\%）。
高信度 Red 阳性率（28.6\%）$>$ Green（5.1\%），证实高信度分区具有临床区分度。

Stage 1B 置信度引擎在架构上吸纳了多定义一致性框架的核心精神：
以四分量（字段完整度、努力程度、锚点一致、结局一致）加权置信度替代 K 定义投票，
提供连续置信度分数而非二元高/低信度标签。

% ============================================================
% 参考文献（独立编译时取消注释）
% ============================================================
% \bibliography{references}

\end{document}
